\section{Introduction}

Light field cameras capture both spatial and angular information of light rays, providing rich geometric cues for 3D reconstruction and depth estimation [1]. Extracting epipolar plane image (EPI)s (EPIs) and analyzing angular signals remain the dominant paradigms for inferring scene geometry. These approaches achieve significant progress in controlled environments where surfaces exhibit ideal diffuse reflection. However, real-world scenes extensively contain non-Lambertian materials, such as glossy, specular, and translucent surfaces. Inferring the accurate depth of these complex materials represents a hard challenge for both traditional algorithms and deep networks [22].

Most current light field depth estimation methods rely heavily on the Lambertian assumption, which requires pixel appearance consistency across different viewpoints [2]. Unfortunately, non-Lambertian surfaces explicitly violate this physical principle. Specular reflections shift dynamically with viewpoint changes, while volumetric scattering prevents the existence of a single surface depth. These physical phenomena destroy the Lambertian angular cone geometric constraints, forming $X$-shaped patterns instead of linear structures in EPIs [3]. Consequently, standard EPI-based architectures suffer from severe depth ambiguity and structural discontinuities when processing reflective or scattering regions.

Beyond the physical assumption violation, non-Lambertian depth estimation faces a severe data scarcity bottleneck. For instance, the widely used Non-Lambertian dataset provides only four training scenes, which falls orders of magnitude below the minimum threshold required for learned generalization [4]. Our extensive preliminary experiments, comprising 132 trials across 16 research directions, quantitatively demonstrate this degradation. While EPI-based methods reduce the mean absolute error ($\text{MAE}$) to $0.081$ in mixed urban scenes, the $\text{MAE}$ substantially increases to $0.411$ in non-Lambertian regions, indicating a performance drop of over five times. Furthermore, complex textures cause EPI slope ambiguity even in Lambertian areas, locking the overall $\text{MAE}$ at $0.387$. These results confirm that the performance degradation stems from inherent architectural limitations and critical data deficiency, rather than suboptimal hyperparameter tuning.

Addressing non-Lambertian light field depth estimation is highly necessary for practical applications like autonomous driving and robotic manipulation, where reflective objects are ubiquitous [5]. Relying on a single EPI branch or blindly adjusting loss functions cannot overcome the violated physical assumptions. Therefore, it is essential to develop a unified framework that decouples light field signals at the physical level, explicitly routes different material types, and maintains robustness under imbalanced data distributions. Establishing such a paradigm not only clarifies the theoretical boundaries of traditional EPI-based methods but also provides a reliable foundation for complex light field scene understanding.

Traditional light field depth estimation algorithms primarily rely on photo-consistency and epipolar plane image (EPI) analysis for inferring scene geometry [6]. Methodologically, these approaches strictly assume Lambertian reflectance, rendering them highly ineffective for processing glossy or scattering surfaces in real scenes. Specifically, specular reflections destroy the linear epipolar structures, forming ambiguous X-shaped patterns that invalidate standard geometric constraints used in matching [23]. Technically, many traditional methods employ frequency domain analysis, such as the 2D discrete Fourier transform, to separate depth and material properties [8]. However, this spectral approach consistently fails under the low angular resolution provided by practical commercial light field camera imaging systems. The discrete 9x9 angular sampling provides insufficient frequency resolution, preventing the accurate decoupling of environmental illumination, scene geometry, and material reflectance components.

Early deep learning methods attempt to address these limitations by extracting features from epipolar images using convolutional neural networks [9]. Methodologically, these data-driven approaches often misinterpret non-Lambertian degradation as a simple network capacity issue or a parameter deficiency problem. Consequently, researchers extensively tune loss functions or increase network depth, largely overlooking the core violation of underlying physical assumptions [10]. Technically, these models predominantly utilize single-branch architectures that process Lambertian and non-Lambertian regions uniformly without any explicit spatial separation mechanism. This unified processing forces the network to learn conflicting feature representations, resulting in severe depth ambiguity at material boundaries [11]. The lack of explicit geometric routing prevents the network from adapting to the distinct X-shaped patterns of highly reflective surfaces.

Recent advanced networks incorporate attention mechanisms, multi-modal inputs, or defocus auxiliary heads to improve overall depth estimation performance significantly [12]. Methodologically, these approaches lack explicit physical decoupling, treating depth estimation as a purely statistical mapping problem without incorporating physical priors. Without separating environmental illumination from material reflectance, these models struggle to generalize effectively across mixed or complex urban scenes [13]. Technically, current methods suffer from severe cross-domain generalization bottlenecks due to highly imbalanced training data distributions across different public datasets. Non-Lambertian datasets contain only a few scenes, causing severe overfitting when models are trained exclusively on specific visual domains [14]. Existing training strategies rarely employ domain-balanced sampling or global statistical validation, limiting their robustness in diverse real-world practical applications.

Therefore, effectively addressing non-Lambertian depth estimation requires much more than incremental architectural tweaks or simple data augmentation techniques alone. It necessitates a unified framework that explicitly decouples physical signals, routes distinct geometric patterns adaptively, and ensures robust cross-domain generalization.

In this paper, we propose a unified dual-mask physical model for light field depth estimation that explicitly decouples environmental illumination, scene geometry, and material reflectance. Instead of relying on ineffective frequency domain analysis under low angular resolution, our approach introduces a three-layer angular signal decomposition to extract robust geometric pseudo-labels. These physical priors guide a dual-branch architecture that adaptively routes Lambertian and non-Lambertian regions to specialized processing streams for targeted feature extraction. By integrating a pixel-wise dual-mask fusion mechanism with a domain-balanced training framework, our model effectively addresses the geometric degradation caused by specular and scattering surfaces in complex 9×9 light field scenes.

\begin{itemize}
  \item \textbf{Contribution 1}: We propose a physics-based three-layer angular signal decomposition model to decouple light field signals into ambient illumination, disparity depth, and material reflectance. Formulated as $I(u,v) = DC + a \cdot u + b \cdot v + \varepsilon(u,v)$, this model extracts geometric features like angular gradients from the residual term. This physical decoupling overcomes the resolution limitations of traditional spectral methods in 9×9 discrete sampling, providing reliable pseudo-labels for accurate material classification and subsequent depth regression.
  \item \textbf{Contribution 2}: We develop a dual-branch adaptive routing architecture, named GeometricDualMask, which addresses the physical failure of epipolar plane image (EPI) constraints on non-Lambertian surfaces. Through extensive empirical analysis, we identify that specular reflections form X-shaped patterns rather than linear structures, violating standard geometric assumptions. Our design leverages bimodal features and X-shaped patterns to route pixels into specialized branches, fusing outputs via a dual-mask mechanism supervised by angular gradients to prevent depth degradation.
  \item \textbf{Contribution 3}: We construct a unified light field dataset loader and a multi-domain balanced training framework to ensure strong cross-domain generalization. By performing global statistical validation on extracted disparity amplitude and material ratio using Cohen's d effect size, we verify the discriminative power of our features across approximately 230 scenes. Furthermore, we implement a domain-balanced sampling strategy with a weighted random sampler to mitigate overfitting caused by scarce non-Lambertian data, enhancing overall model robustness.
\end{itemize}
Extensive experiments on multiple cross-domain datasets demonstrate that our proposed method achieves an overall Mean Absolute Error (MAE) of 0.133, significantly outperforming existing baselines. Specifically, the model yields an MAE of less than 0.16 for Lambertian scenes and maintains robust performance with an MAE below 0.25 for challenging non-Lambertian regions. These quantitative results validate the effectiveness of our unified physical model and dual-branch routing mechanism in handling complex real-world light field scenarios. To contextualize these advancements, we begin by reviewing the foundational and traditional approaches that have historically guided light field depth estimation.